\section{Reducing the cost of the exact stage}
\label{sec:fast-method}
The preceding comparison shows that an exact formula can save algebraic work while still spending time on general symbolic setup. We therefore evaluate the finite products in the derived tadpole, one-mass vacuum, and massless propagator formulas directly. Their coefficients remain rational functions of the dimension and the permitted scales. The formulas include the supported quadratic numerators described in Appendix~\ref{app:blocks13}. Other numerator structures and coupled families pass to the existing reduction methods.

For scalar one-loop blocks, we also combine repeated terms before applying the recurrence in Eq.~\eqref{eq:analytic-recurrence}. The coefficients are polynomials in $d$ with exact rational coefficients, so their arithmetic can use integer numerators and denominators throughout. Appendix~\ref{app:fast-method} derives the combined recurrence. Independent blocks share this method, and repeated identical blocks share intermediate results within one request. No reduction coefficient from an earlier request enters a new answer.

Several raised indices can supply the next recurrence step. Each choice gives a valid decomposition, but the number of intermediate integrals can differ. We train a cost-sensitive tree to select among four strategies: choose the largest positive power, the smallest power greater than one, the first raised index, or the last raised index. A strategy applies at each intermediate integral. The model selects one strategy for the target collection before the exact recurrence begins; it does not predict a coefficient. The eleven inputs describe the number of denominators and targets, the extra positive powers, and the powers at the first five positions. Appendix~\ref{app:fast-method} gives the features and training objective.

We measure all four strategies on 112 training requests and 56 separate validation requests. Each request contains one or two scalar targets with two through five denominators and at most six extra positive powers. Three repetitions provide a median recurrence time, and all strategies must agree exactly. Validation selects the tree and a minimum predicted saving, with the best fixed strategy as its control. These labels measure recurrence time plus prediction overhead. The reserved application comparison then tests whether the selected policy reduces total time after loading, basis checks, and output verification.

The resulting application first splits independent loop blocks, then attempts a direct formula or the scalar recurrence. It loads the existing learned Kira application when a block needs the remaining exact or native methods. This order reduces setup work on inputs with an immediate analytic decomposition. The retained sparse-seed model keeps its earlier weights. Learned row selection remains omitted, and the equation filter retains its separately tested application. The new recurrence policy adds a different decision to this combination.

Complete application validation also governs deployment of the recurrence model. Section~\ref{sec:fast-deployment} measures that decision on the original validation split, freezes its outcome, and tests twelve fresh requests.
